\documentclass[11pt]{article}

\usepackage{amsmath}
\usepackage{amssymb}
\usepackage{geometry}
\usepackage{graphicx}
\usepackage{hyperref}
\usepackage{listings}
\usepackage{hyperref}

\title{Trabajo práctico Organización de Computadoras}
\author{Juan Manuel Antuña \and Isaías Cheij \and Tomás Miranda \and Hilario Yáñez Bertola}
\date{Junio de 2026}

\begin{document}
    \maketitle
    \newpage

    \section{Primera parte: Un algoritmo de multiplicación microprogramado}

    En esta primera parte vamos a desarrollar un (o un par de) algoritmo(s) de multiplicación basados en microinstrucciones.
    Puntualmente, vamos a implementar una multiplicación para números enteros de 8 bits, representados en complemento a la base (two's complement).
    Usamos como base el procesador desarrollado durante la materia en \href{https://www.cs.colby.edu/djskrien/CPUSim/}{CPUSim}
\end{document}
